% ПРИЛОЖЕНИЯ, без номер на раздел.
% Изходният текст се вмъква с \lstinputlisting от src/ и include/, никога чрез
% копиране: копието се разминава с оригинала при първата поправка.
\section*{ПРИЛОЖЕНИЯ}
\label{sec:appendix}
\addcontentsline{toc}{section}{ПРИЛОЖЕНИЯ}

\subsection*{Приложение А. Договори на рамката}

Двата договора, върху които е построено всичко останало: \code{Problem} описва задача,
\code{Metaheuristic} описва търсене. Никой от двата не знае нищо за нишки.

\lstinputlisting[caption={\code{include/anneal/core.hpp}}, label={lst:core}]%
{../include/anneal/core.hpp}

\subsection*{Приложение Б. Споделено най-добро решение}

Единственото споделено изменяемо състояние в кооперативната схема. Атомарната цена е
бързият път, а решението се записва под мютекс с повторна проверка вътре в критичната
секция; мотивът е обяснен в §\ref{subsec:impl-threads}.

\lstinputlisting[caption={\code{src/parallel.cpp}, \code{SharedBest::offer} и \code{SharedBest::read}},
                 label={lst:sharedbest}, firstline=54, lastline=105, firstnumber=54]%
{../src/parallel.cpp}

\subsection*{Приложение В. Обща част на трите схеми}

Трите схеми се различават само по обратното извикване \code{exchange}. Всичко останало,
тоест създаването на работниците, спирането по бюджет и събирането на резултата, е общо
и е написано веднъж.

\lstinputlisting[caption={\code{src/parallel.cpp}, \code{run\_workers}}, label={lst:runworkers},
                 firstline=137, lastline=213, firstnumber=137]%
{../src/parallel.cpp}

\subsection*{Приложение Г. Островен модел с миграция по пръстен}

\lstinputlisting[caption={\code{src/parallel.cpp}, \code{run\_island}}, label={lst:island},
                 firstline=260, lastline=310, firstnumber=260]%
{../src/parallel.cpp}

\subsection*{Приложение Д. Генератор на съседен ход при търговския пътник}

Функцията с най-голямо измерено влияние върху качеството, заедно с коментара, който
записва защо първата ѝ версия беше грешна.

\lstinputlisting[caption={\code{src/tsp.cpp}, \code{random\_move} и \code{delta}}, label={lst:move},
                 firstline=160, lastline=236, firstnumber=160]%
{../src/tsp.cpp}

\subsection*{Приложение Е. Използвани еталонни инстанции}

Всички инстанции се генерират от програмата. Произходът на обявената стойност за всяка
от тях е даден в Табл.~\ref{tab:instances}. Реализиран е и четец на формата
TSPLIB~\cite{reinelt1991tsplib} за подмножеството EUC\_2D, за да може публикувана
инстанция да бъде подадена отвън, но нито едно измерване в Раздел~\ref{sec:results} не
зависи от изтеглен файл.

\subsection*{Приложение Ж. Команди за възпроизвеждане на измерванията}

Всички команди са изпълнявани от главната директория на хранилището. Версиите на
инструментите са в Табл.~\ref{tab:environment}.

\begin{lstlisting}[language=bash, caption={Изграждане, тестове, демонстрация}, label={lst:build}]
cmake -S . -B build -G Ninja -DCMAKE_BUILD_TYPE=Release
cmake --build build
ctest --test-dir build --output-on-failure
cmake --build build --target demo
\end{lstlisting}

\begin{lstlisting}[language=bash, caption={Измервания, по едно за всяка таблица}, label={lst:bench}]
./build/anneal_bench objective       # Табл. 6.1 и 6.2
./build/anneal_bench neighbourhood   # Табл. 6.3
./build/anneal_bench quality         # Табл. 6.4
./build/anneal_bench schemes         # Табл. 6.5
./build/anneal_bench speedup         # Табл. 6.6 и Фиг. 6.1
./build/anneal_bench falsesharing    # Табл. 6.7
./build/anneal_bench profile         # Табл. 6.8
./build/anneal_bench all             # всички наведнъж
\end{lstlisting}

\subsection*{Приложение З. Изход на демонстрационната програма}

Изходът по-долу е дословен от изпълнение на машината от Табл.~\ref{tab:environment}.
Демонстрацията избира сама най-добрия по медиана алгоритъм в част 1 и го използва в
останалите части, тоест изборът не е фиксиран в текста.

\lstinputlisting[caption={\code{./build/anneal\_demo}}, label={lst:demo},
                 basicstyle=\tiny\ttfamily, numbers=none, language={}]%
{demo_output.txt}

\subsection*{Приложение И. Изпълнител на тестовете}

Изпълнителят е деветдесет реда и няма външна зависимост. Това е съзнателно: рамка за
тестове, изтегляна при конфигуриране, е зависимост, която трето лице трябва да
удовлетвори, преди първото изграждане изобщо да проработи.

\lstinputlisting[caption={\code{tests/check.hpp}}, label={lst:check}]%
{../tests/check.hpp}
